\section{Catalog and Replication}
\label{sec:catalog}

\subsection{Records}

A \emph{store record} contains a display name, a region, two direct-payment strings (a BOLT 12 offer and a BIP 352 code), a content rating for the whole store, and the store's 32-byte x-only public key $S$, which is the store identifier. Version 3 adds the mailbox key of \S\ref{sec:identity}, the store's onion host, the moderators the store accepts (by their signing keys), a public profile (tagline, about text, avatar and banner named by the SHA-256 of their bytes, shipping and returns policies, processing time, the countries it ships to, and the day it was created), and one byte of payment-rail flags (\S\ref{sec:payments}). A \emph{listing} contains $S$, a title, a category, a price in satoshis, a ship-from region, a description, a seller display name, a stock count, the date it was signed, whether it accepts a moderator, a kind (physical or digital), an optional shared product key, flat shipping in satoshis, the subtotal above which shipping is waived, and the vendor's content label. Neither record has a field for a legal name, an email address, a phone number, or a street address, and signing refuses free-text fields that contain them.

\paragraph{Plain-text hardening.} Because records are replicated across nodes and rendered in other users' interfaces, every public text field---across store records, listings, version-7 commerce fields (condition notes, specifics, variants, shipping services, pickup areas, and return policies), and seller or ship-from filter values---must pass a \code{plain\_text} check refusing ASCII and C1 control characters (\code{U+0000}--\code{U+001F}, \code{U+007F}--\code{U+009F}) and Unicode bidirectional override and isolate characters (\code{U+061C}, \code{U+200E}--\code{U+200F}, \code{U+2028}--\code{U+202E}, \code{U+2066}--\code{U+2069}); multiline descriptions may keep tabs and line breaks, and ordinary Unicode including emoji (zero-width joiners and variation selectors) passes. Unlike the heuristic contact-and-location check, which runs only at signing, the \code{plain\_text} check runs on \emph{both signing and verification}, so nodes and browsers drop hostile records rather than displaying or relaying them.

\subsection{Canonical encoding}

Each record body is a fixed sequence of fields: a four-byte magic (\code{OBST} for stores, \code{OBLS} for listings, \code{OBLB} for label sets, \code{OBCX} for catalog indexes, \code{OBAX} for admission snapshots), a one-byte version, then fields in a fixed order. Integers are little-endian of fixed width. Strings are a two-byte little-endian length followed by UTF-8 bytes; blobs are a four-byte length, capped at 8~MiB. The decoder rejects a wrong magic, an unknown version, invalid UTF-8, truncation, and trailing bytes. Because every variable-length field is length-prefixed and the field order is fixed, distinct records have distinct encodings, and a decoded record re-encodes to the same bytes. Appendix~\ref{app:formats} gives the layouts. A new version adds fields at the end of the body; older versions remain decodable, so records signed before a field existed keep verifying.

\subsection{Signatures and addresses}

The signature is a BIP 340 Schnorr signature~\citep{bip340} over a tagged hash of the body:
\[
  \sigma = \mathsf{Sign}_{s}\bigl(\Hash_{\mathrm{tag}}(\mathit{body})\bigr),\quad \Hash_{\mathrm{tag}}(m)=\mathrm{SHA256}\bigl(\mathrm{SHA256}(\mathrm{tag})\,\|\,\mathrm{SHA256}(\mathrm{tag})\,\|\,m\bigr),
\]
with tags \code{ob/store/v1}, \code{ob/listing/v1}, and \code{ob/labels/v1}, so a signature over one record type is not valid for another. Signing uses BIP 340's deterministic nonce derivation without auxiliary randomness. The signature is appended to the body. A listing can only be signed by the secret key of the store it names, and verification checks the signature against the $S$ inside the listing body.

A listing's address is $\mathrm{SHA256}(\mathit{body})$, computed over the body without the signature. An order commits to this address, so a vendor who later edits a listing (for example, raising the price) produces a new address and cannot claim that the buyer agreed to the new terms. A node refuses an order for an address that it has not verified.

\subsection{Content labels and labelers}
\label{sec:labels}

The network carries ordinary goods, but ``ordinary'' differs by buyer: some do not want to see weapons, alcohol, or adult material, and some jurisdictions restrict them. The design answers this with labels and filtering on the buyer's own device, not with a central moderator of the catalog.

\paragraph{Self-labels.} Each listing carries a content rating (\emph{general}, \emph{sensitive}, or \emph{adult}) and a 16-bit set of tags (nudity, violence, drugs, weapons, gambling; the other bits are reserved). The store record carries a rating that applies to everything the store sells. Both are inside the signed body, so a replica cannot strip or soften a label without breaking the signature, and a vendor who wants a different label must sign a new record (for a listing, at a new address).

\paragraph{Labelers.} A \emph{labeler} is a BIP 340 key that publishes numbered \emph{label sets}. Each entry in a set names a target (a listing address or a store key) and either gives it a rating, a tag set, and a short reason code, or retracts the labeler's earlier label on that target. For each labeler and target, the entry in the highest-numbered set that names the target is the current one, so a later set may carry only changes, and a retraction leaves no label. Entries are sorted by target with no repeats, so one set has one encoding; a set is addressed by the hash of its body like a listing. A reason is a short code from a fixed alphabet, not free text, so a label set cannot be used to publish arbitrary statements about a vendor.

\paragraph{Filtering.} A buyer chooses which labelers to follow; the follow list is private to the buyer's node. The client combines the store rating, the listing's self-label, the labels from followed labelers, and the buyer's own blocked words and hidden stores and listings, and decides to show, cover (show a placeholder the buyer can open), or hide each listing. When several rules apply, the strictest wins, and the client shows which source caused it. Safe mode is on by default and hides sensitive and adult listings; the buyer can relax each rule.

\paragraph{What labels do not do.} A label changes only what one device shows. It does not remove a record from any replica, block an order, or bind other buyers. A vendor can mislabel its own goods; a labeler can label wrongly or selectively; and a buyer who follows no labeler relies only on self-labels and its own lists. The design accepts this: it gives buyers tools to filter without giving any party the power to delete. Fetching a labeler's sets reveals to whoever serves them that the node follows that labeler, so sets should be fetched through the same transport as catalog records.

\subsection{Signed catalog indexes and freshness}
\label{sec:catalog-index}

A listing signature alone proves only that the store key holder signed those bytes at some time; without an index, a replica can keep serving retracted listings or omit current ones without failing individual signature checks. To close this gap, each store signs a \emph{catalog index} (\code{OBCX} version 1, tag \code{ob/catalog-index/v1}): the store identifier $S$, a monotonically increasing sequence number, a signed-on day, the SHA-256 digest of the store's current encoded store record, and the sorted, deduplicated list of up to 4{,}096 current listing content hashes. It carries no free text.

Nodes and read-only browsers keep the highest verifying sequence for each store and verify the index's signature against $S$. When a catalog index matches the store record, any listing served by a replica that is absent from the index---or that fails signature verification---is dropped, so retracted listings disappear and a replica cannot splice older versions into a current storefront. What remains outside the signature is global freshness: a replica or colluding introductions can still withhold a newer index sequence, though querying the store's own onion and rotating across multiple introductions raises the bar to withholding across all sources (\S\ref{sec:limitations}).

\subsection{Replication}

Every node stores the signed headers of the catalog it has seen, so that a store remains browsable while its publisher is offline, and images and files must be held by at least one replica before a listing counts as published. Because records are self-authenticating, a replica needs no trust for integrity: a client or node that receives a record from anyone checks the signature, the content hash, the plain-text rules, and the store's catalog index when present. Replicas are trusted only for availability and see who fetches what (\S\ref{sec:privacy}).

An open catalog invites spam. The directory below answers this with a fee that every node checks on chain. A node replicates store records for stores admitted to the directory and stores it follows, and listings only for its own store, the stores it follows, and stores recently opened on it.

\subsection{The directory}
\label{sec:directory}

Following a store needs its address, which a buyer must learn somewhere. The directory is how a node learns stores it has not been told about, without a server that can alter what they publish.

\paragraph{The head and admission snapshot.} A directory is named by one BIP 340 key, the \emph{document key}, which ships with the software. Its operator signs a \emph{head} (\code{OBSD} version 3, tag \code{ob/seed/v3}): a sequence number, a sorted list of up to 64 introduction onions, and a schedule of up to eight fee terms entries. Each entry names a fee key $K_F$, an amount, a confirmation depth, an epoch in blocks (at most 52{,}416, about a year), the height it starts at, and the height it was announced at. The first head names 20{,}000 satoshis, one confirmation, and an epoch of 2{,}016 blocks, about two weeks. A head names no stores. Alongside the head, the directory key signs a companion \emph{admission snapshot} (\code{OBAX} version 1, tag \code{ob/admission-index/v1}): the directory signer key, a sequence number, a signed-on day, the head sequence it accompanies, and the sorted, deduplicated list of up to 4{,}096 admitted store identifiers. A node keeps the highest head sequence and highest admission snapshot sequence that verify under the document key. A later head replaces an earlier one only if it extends the schedule the node holds: it keeps every held entry, may drop leading entries under which no store can still be listed, and adds entries only at the end. An added entry must start at least 1{,}008 blocks after its announcement if it asks more of a store (a new key, a higher fee, more confirmations, or a shorter epoch) and 288 blocks otherwise, and a node refuses one announced more than 144 blocks before the tip at which it last saw the head current. Two heads with one sequence are refused and kept as evidence. Two introduction onions ship with the software for a node that holds no head yet; after that, the node dials the list in the head it holds.

\paragraph{Introductions.} An \emph{introduction} is an ordinary node whose onion the head lists. It stores and serves bytes that verify under their own keys (the head, the admission snapshot, admissions, store records, catalog indexes, and listings) and signs none of them. Periodically a node asks up to three introductions, rotating through the list, for a newer head, the signed admission snapshot, pages of admissions ordered by store key, the store record of every admitted or followed store it lacks, and catalog indexes for those stores. It fetches listings only for its own store, the stores it follows, and the 50 stores most recently opened on it; opening a store's page fetches its catalog index and listings from the onion in the store's signed record, and from an introduction if that onion does not answer. A seller's node submits its store record, its admission, its signed catalog index, and its listings to the introductions it asks. A source that does not answer is skipped, and a record that does not verify never replaces one that did. All of this runs over Tor. A node can turn it off; it then learns no stores from the directory and is not listed there, and a store followed by its address still works.

\paragraph{Admission and the fee.} A store joins by paying the fee, not by asking the operator. Each store has its own fee address: a key-path-only Taproot output for
\[
  K_{F,S} = K_F + t\,G,\qquad t = \Hash_{\mathrm{ob/fee\text{-}key/v1}}\bigl(K_F \,\|\, S\bigr),
\]
with $K_F$ and $S$ in x-only form. Anyone can compute $K_{F,S}$ from the entry and the store key. Only the holder of the fee key's secret $k_F$ can spend it, with the secret $k_F+t$ (negated as BIP 340 requires for an x-only key). Because the address names the store, the fee can be paid from any wallet, exchange, or swap service in one ordinary transaction with no marker output, and one payment cannot admit a different store. Once the payment has the entry's confirmation depth, the store key signs an \emph{admission} (\code{OBAD} version 2, tag \code{ob/admit/v2}): the store key, the transaction id, the output index, and the height of the block that confirmed it. The seller's node pays from its own wallet when it can, accepts a payment to the store's fee address from anywhere, and pays again when the epoch runs out. Within six blocks of a scheduled entry it pays that entry's terms if they are higher. An earlier encoding, admission version 1, pays $K_F$ itself and names the store in an \code{OP\_RETURN} output carrying \code{OBAD} and $S$; nodes still accept it.

\paragraph{Every node checks the chain.} An introduction keeps any admission whose signature verifies. Each node then decides for itself. An admission is judged by the schedule entry in force at the height that confirmed its payment, not by the newest entry. It counts only if its transaction is confirmed at the stated height, that height is within that entry's epoch of the node's own chain tip and has at least its confirmation depth, and the named output pays at least its fee to the store's fee address (for version 1, to $K_F$, with the store's \code{OP\_RETURN} in the same transaction). A node shows a store when it holds a verified store record and either follows the store or holds an admission that passes this check. The signed admission snapshot (\code{OBAX}) helps nodes and read-only browsers learn which stores are admitted and audit omissions, but it is advisory for chain validity: an admission whose chain check fails is still dropped regardless of what the snapshot names. When the epoch passes, the admission lapses until the store pays again. The fee makes publishing at scale cost money; it is not a bond and says nothing about whether a seller is honest.

\paragraph{What the operator controls.} The operator sets the fee terms, the introduction list, and the admission snapshot, and receives the fees. It cannot alter or forge a store record, catalog index, listing, or admission, since each verifies under its own key, and it cannot force nodes to accept an unpaid admission because each node checks the chain. It can change the fee terms only for every store at once and only going forward: a new entry applies to payments confirmed after it starts, with notice, and a store that already paid stays listed for the epoch it paid for. The fee key cannot change yet; the software refuses to sign an entry with a new one. What the operator and the introductions can do is withhold. An introduction can decline to serve a store's records, as any replica can; a store whose records every listed introduction withholds is not found through the directory, while a buyer who has its address can still open it (\S\ref{sec:discovery}). Because the directory signs \code{OBAX}, a store with a valid on-chain admission that is absent from the snapshot leaves signed, checkable evidence of exclusion, whereas before \code{OBAX} omission was silent. A node accepts heads and snapshots only under the document key it was built with, so a release has one directory (\S\ref{sec:limitations}).
